\documentclass[a4paper]{amsart}
\usepackage{subfiles}
\usepackage[margin=2cm]{geometry} % Adjust margin as needed

\usepackage{todonotes}
\usepackage{amsmath}
\usepackage{amssymb}
\usepackage{amsthm}
\newtheorem{theorem}{Theorem}

\newcommand\meet{\mathrm{meet}}
\newcommand\todoin[2][]{\todo[inline, caption={2do}, #1]{
        \begin{minipage}{\textwidth-4pt}#2\end{minipage}}}


\begin{document}

The meet computation in typclingo.

Let $V$ be a set of \emph{type variables},
$F$ be a set of \emph{function constructors}, and
$T$ be a set of \emph{type constants}.

The set $\mathcal{T}$ of \emph{type expressions} is the smallest set containing:
\begin{itemize}
	\item $\top$ (the top type)
	\item $e$ for $e \in T \cup V$
	\item $f(e_1,\dots,e_n)$ for $e_1,\dots,e_n \in \mathcal{T}$, $f \in F$, and $n \geq 0$ (function types)
	\item $(e_1 \mid ... \mid e_n)$ for $e_1,\dots,e_n \in \mathcal{T}$ and $n \geq 0$ (type union)
\end{itemize}

We omit the parentheses for function types if the arity is zero
and type unions when there is no risk of ambiguity.
We define $\bot$ as the empty type union.

A \emph{type specification} $S$ is a finite set of expressions $t \lessdot e$ or $t \doteq e$
such that $t \in T$ and $e \in \mathcal{T}$ where $e$ does not contain type variables.

We assume that the type constants always contain $\mathbb{F}$, $\mathbb{N}$ and $\mathbb{S}$.
Furthermore, type specifications always include:
\begin{itemize}
	\item $\mathbb{F} \lessdot \top$ (the type containing all function types)
	\item $\mathbb{N} \lessdot \top$ (the type containing all numbers)
	\item $\mathbb{S} \doteq \mathbb{F} \mid \mathbb{N}$
\end{itemize}

Each type constant occurring on the right-hand side of a $\lessdot$ or $\doteq$ is defined exactly once.

We let $\Sigma$ be the set of all \emph{type variable mappings} from type variables to type expressions.
We assume that mappings $\sigma\in\Sigma$ map all variables in $V$ to some type expression.
We denote a type mapping as a set of expressions of the form $v \mapsto e$ where $v \in V$ and $e \in \mathcal{T}$.
A mapping maps type variables to $\mathbb{S}$ if not explicitly given.
Given $\sigma_1, \sigma_2 \in \Sigma$, we let $\sigma_1 \sqcup \sigma_2$ be the mapping defined as:
$(\sigma_1 \sqcup \sigma_2)(v) = \sigma_1(v) \mid \sigma_2(v)$ for all $v \in V$.

To handle recursive types,
we use a call guard set $C \subseteq T \times \mathcal{T}$ of pairs of type constants and type expressions
for which the meet is currently being computed.
Function $\mathrm{meet}_{S}: \mathcal{T} \times \mathcal{T} \times \Sigma \times C \mapsto \mathcal{T} \times \Sigma$
computes the \emph{meet} of two type expressions given a type variable mapping, a type specification $S$, and a call guard set $C$.
The result is a new type expression and an updated type variable mapping.

We define $\mathrm{meet}_{S}$ as follows:
\begin{itemize}
	\item $\meet_{S}(e, \top, \sigma, C) = \meet_{S}(\top, e, \sigma, C) = (e, \sigma)$
	\item $\meet_S(t, e_1, \sigma, C) = \meet_S(e_1, t, \sigma, C) = \meet_S(e_1,e_2,\sigma,C\cup\{(t,e_1)\})$\\
	      where $t \doteq e_2 \in S$ and $(t,e_1) \notin C$
	\item $\meet_{S}(v, e, \sigma, C) = \meet_{S}(e, v, \sigma, C) = (e', \{v' \mapsto \sigma'(v') \mid v' \neq v \in V\} \cup \{v \mapsto e'\})$\\
	      where $v \in V$, and $(e', \sigma') = \meet_{S}(\sigma(v), e, \sigma, C)$
	\item $\meet_S(e, e_1 \mid \dots \mid e_n, \sigma, C) = \meet_S(e_1 \mid \dots \mid e_n, e, \sigma, C) = (e_{i_1}' \mid \dots \mid e_{i_k}', \sigma_{i_1} \sqcup \dots \sqcup \sigma_{i_k})$ \\
	      where $(e_i',\sigma_i) = \meet_S(e_i,e,\sigma,C)$ and $i_j$ are all the indices such that $e_{i_j}' \neq \bot$
	\item $\meet_S(f(e_1,\dots,e_n),\mathbb{F},\sigma,C) = \meet_S(\mathbb{F}, f(e_1,\dots,e_n),\sigma,C) = (f(e_1,\dots,e_n), \sigma)$
	\item $\meet_S(f(e_1,\dots,e_n),t,\sigma,C) = \meet_S(t,f(e_1,\dots,e_n),\sigma,C) = \begin{cases} (t,\sigma') & e' \neq \bot \\ (\bot,\sigma') & \text{otherwise} \end{cases}$\\
	      where $t \lessdot e \in S$, $t \notin \{\mathbb{N}, \mathbb{F}\}$, $(t,f(e_1,\dots,e_n)) \notin C$, and $(e',\sigma') = \meet_S(f(e_1,\dots,e_n),e,\sigma,C \cup \{(t,f(e_1,\dots,e_n))\})$
	\item $\meet_S(f(e_1,\dots,e_n),f(g_1,\dots,g_n), \sigma, C) = \begin{cases} (f(h_1,\dots,h_n),\sigma_n) & \text{all } h_i \neq \bot \\ (\bot,\sigma_n) & \text{otherwise} \end{cases}$\\
	      where $\sigma_0 = \sigma$ and $(h_i, \sigma_i) = \meet_S(e_i,g_i,\sigma_{i-1}, C)$
	\item $\meet_S(t_1,t_2,\sigma,C) = \begin{cases} (t_2, \sigma) & t_1 \text{ is a reachable super type of } t_2 \text{ in } S \\ (t_1, \sigma) & t_2 \text{ is a reachable super type of } t_1 \text{ in } S \\ (\bot,\sigma) & \text{otherwise} \end{cases}$\\
	      where $t_1 \lessdot e_1, t_2 \lessdot e_2 \in S$
	\item $\meet_S(e_1,e_2,\sigma,C) = (\bot,\sigma)$, otherwise
\end{itemize}

To give an idea how the meet computation is applied in typclingo,
let $p(X), q(X)$ be some rule body.
We assume that we have some type specification $S$ containing type definitions $e_1$ and $e_2$
for the arguments of $p$ and $q$, respectively.
We can then map the atoms $p(X)$ and $q(X)$ to the type expressions $p(x)$ and $q(x)$ where $x$ is a type variable.

To check that the body is well-typed,
we compute $\mathrm{meet}_{S}(p(x), p(e_1),\emptyset,\emptyset) = (e,\sigma)$.
In the next step, we compute $\mathrm{meet}_{S}(q(x), q(e_2),\sigma,\emptyset) = (e',\sigma')$.
If both $e$ and $e'$ are not $\bot$,
then the body is well-typed and the type of $X$ is given by $\sigma'(x)$.

\newpage
\subfile{proof-termination}
\end{document}
